\section{Discussion}
\label{sec:discussion}

{\bf Does ``reads LLM-written'' change safety-relevant behaviour?}
For \llama, \qwen and \luna it does not, at the resolution of this study (CIs of about $\pm$2--5\pp, 147--350 items per cell, two rewriters).
For \gemini it does, slightly: refusal rises by 3--4\pp on both harmful and benign requests.
This holds even though \luna recognises LLM style best of all four models, so recognising the style is not enough to change behaviour.

{\bf Is the change carried by an authorship representation?}
Four lines of evidence say no.
First, a deterministic polite opener on the original text reproduces \gemini's over-refusal effect.
Second, \gemini does not judge those prompts AI-written (1.5--6\%), and an explicit ``an AI agent wrote this'' label does nothing.
Third, across all four models the explicit label, the purest manipulation of belief about the author, never produces a Holm-significant effect.
Fourth, in the open models the direction that tracks ``reads LLM'' is a register axis: it fails to detect LLM-authored text written in human style, steering it at natural magnitudes does little beyond random directions, and in \qwen steering flips the model's authorship judgement without changing behaviour.
The strong effect it does have, in \llama toward casual register, works through register mirroring rather than authorship belief.

{\bf Relation to evaluation awareness.}
The style direction is nearly orthogonal to the evaluation-awareness direction in both models, so LLM-written-looking prompts are not processed as ``test-like'' along that axis.
This agrees with \citet{devbunova2026evaluation}, whose evaluation probe tracked benchmark format (multiple-choice structure) rather than register; our \vL and \vS variants are formal but not benchmark-formatted.
More broadly, our results fit the pattern that apparent ``source'' effects reduce to surface features~\citep{chalkidis2026templated,ackerman2024inspection,devbunova2026evaluation}, and extend it from judged items to the instruction itself and to four behaviours.
Like \citet{xu2025ai}, we see more style sensitivity in a larger API model (\gemini) than in 7--8B models, but not in \luna.

\para{Implications for evaluation practice.}
Safety evaluations that use LLM-written user turns should transfer to human users to within a few percentage points for the models we test.
Two caveats matter in practice.
Model-specific sensitivity to politeness and request framing can shift refusal by a few points, so evaluators should match the \emph{register} of deployment traffic, not only its content.
Deployed safety stacks are also style-sensitive: \luna's provider filter blocked human-style harmful prompts about 1.7$\times$ as often as LLM-style ones, so evaluations with LLM-written prompts would under-count filter refusals.
Finally, the human-ward steering result shows that a casual-persona direction can act as a strong jailbreak in \llama when pushed beyond natural magnitudes; we report it because it is relevant to defenders who monitor activation-level attacks.

\subsection{Limitations}
\label{sec:limitations}

\para{Effect-size resolution.}
With 150--350 items per cell we can rule out average effects larger than about 4--5\pp; smaller effects, or effects confined to sub-categories, remain possible.
Harmful refusal sits near ceiling (86--95\%), which limits headroom.

\para{Operationalisation of ``LLM style''.}
Our \vL and \vS variants capture one notion of LLM-written prompts: polite, explicit and formal, from two 2026 rewriters.
Model-written evaluations, structured orchestrator messages with headers and role framing, or self-generated prompts may differ; we excluded headers and mentions of AI or testing to keep content fixed, so the structured agent-message format is untested.

\para{Bundled features.}
Casing and politeness are collinear with style in the rewrites, so only \vPOL, the labels and steering separate them.
\vPOL itself raises \luna's perceived LLM-ness (38.8\% vs.\ 20.6\% for \vORIG), so it is not a pure non-authorship control; the argument for \gemini rests on \gemini neither perceiving \vPOL as AI-written nor reacting to the label.

\para{Judges.}
A single LLM judge, blind to condition, was validated only against \jbb harmfulness labels.
Within-item contrasts are robust to a constant judge bias but not to style-dependent judging, because the judge sees responses whose style may mirror the prompt.

\para{Steering and models.}
We use mean-difference directions in one band of layers with additive steering; doses were calibrated for coherence, not tuned for effect, and the random baseline uses only 2--3 seeds per dose.
Steering was possible only for the two 7--8B open models.
\gemini's reasoning could not be disabled, and 37 of its responses (0.3\%) were empty provider errors.
Decoding was deterministic, so variance comes from items and rewriters, not from sampling.
